% NeurIPS 2026 submission. Requires neurips_2026.sty alongside this file.
% Compile with: pdflatex paper.tex   (run twice for refs)
\documentclass{article}

\usepackage[preprint,nonatbib]{neurips_2026}  % preprint for arXiv; nonatbib lets us use a plain numerical bibliography

\usepackage[utf8]{inputenc}
\usepackage[T1]{fontenc}
\usepackage{microtype}
\usepackage{graphicx}
\usepackage{booktabs}
\usepackage{amsmath,amssymb}
\usepackage{xcolor}
\usepackage{hyperref}
\usepackage{algorithm}
\usepackage{algpseudocode}
\usepackage{url}
\usepackage{mathtools}
\usepackage{placeins}  % \FloatBarrier — keeps figures from drifting across sections
\usepackage{subcaption}

% Macros
\newcommand{\method}{QAttention}
\newcommand{\methodshort}{QAttention}
\newcommand{\bigO}[1]{\mathcal{O}(#1)}
\newcommand{\numtokens}{N}
\newcommand{\depth}{d}
\newcommand{\branching}{b}
\newcommand{\headdim}{D}
\newcommand{\nheads}{H}
\newcommand{\batchsize}{B}
\newcommand{\prefixlen}{L}
\DeclareMathOperator*{\softmax}{softmax}
\DeclareMathOperator{\parent}{parent}

\graphicspath{{figures/}}

\title{QAttention: Linear Tree-Sparse Attention\\
       and Acceptance Decay in Speculative Decoding}

\author{
  Seifeldin Abdellatif\\
  Al Ain University\\
  \texttt{contact@saifmb.com}
  }

\begin{document}
\maketitle

\begin{abstract}
Tree-based speculative decoding accelerates LLM inference by verifying
draft tokens in parallel. Attention doesn't always dominate verifier cost
--- but in practice, the kernels that production systems use
(FlashInfer, DeFT) scale \emph{quadratically} in tree size, so attention
is precisely what blocks the deeper trees that would amortise the static
cost of the target-model forward pass.
We present \textbf{\method{}}, an ancestor-sparse Triton kernel that
scales as $\bigO{\numtokens \!\cdot\! \depth}$ instead of
$\bigO{\numtokens^2}$. \method{} delivers up to
$\boldsymbol{7.75\!\times}$ kernel speedup over FlashInfer and
$\boldsymbol{16.99\!\times}$ over DeFT in the deep/wide regime.
With attention removed as the bottleneck, the end-to-end picture
sharpens: at short context, attention is a small fraction of step time
and the end-to-end gain is a modest $\approx\!1.04\times$ (Amdahl-limited),
but as context grows the share rises and \method{} delivers up to
$\boldsymbol{1.61\!\times}$ end-to-end on a $\branching{=}20,\depth{=}20$
EAGLE-3 workload at $\prefixlen{=}10\,$k, and
$\boldsymbol{1.40\!\times}$ on Sequoia at tree size $1024$. This removes
the systems barrier that has kept tree optimisers in a shallow-tree regime.
The remaining ceiling is acceptance-rate decay: under current draft
models accepted tokens per step plateau as trees grow. We argue this
makes draft-model fidelity, not attention, the next frontier.
\end{abstract}

\section{Introduction}
\label{sec:intro}

Speculative decoding~\cite{leviathan2023fast,chen2023accelerating}
reduces the latency of autoregressive generation by drafting candidate
tokens with a small model and \emph{verifying} them in parallel with the
large target model. Modern systems organise candidates into a
\emph{tree}~\cite{cai2024medusa,li2024eagle,li2024eagle2,li2025eagle3,
chen2024sequoia,miao2024specinfer}: a tree of $\numtokens$ nodes lets the
verifier consider $\numtokens$ continuations in a single forward pass.

\paragraph{Attention is the systematic blocker.}
Attention is not, in principle, the dominant cost of a verifier step.
But the kernels that production stacks rely on for tree
verification --- FlashInfer~\cite{ye2025flashinfer} and
DeFT~\cite{yao2024deft} --- scale as $\bigO{\numtokens^2}$: they iterate
over all KV positions and apply a tree mask. As $\numtokens$ grows, the
attention term dominates super-linearly. Deeper trees are precisely
what is needed to amortise the (essentially static) target-model
forward pass over more accepted tokens; a quadratic-in-$\numtokens$
attention cost makes that amortisation arithmetically impossible.
Empirically the field has settled on small, shallow trees not because
deep trees are algorithmically inferior but because the kernels make
them prohibitively expensive.

\paragraph{Ancestor sparsity.}
In a BFS-ordered $\branching$-ary tree, query at position $q$ may only
attend to the $\depth{+}1$ positions on its ancestor chain
($\parent(k){=}\lfloor (k{-}1)/\branching\rfloor$). The useful work is
$\bigO{\numtokens \!\cdot\! \depth}$, not $\bigO{\numtokens^2}$.
\method{} replaces FlashInfer's KV-block sweep with an explicit
ancestor walk inside an otherwise-standard FlashAttention-2~\cite{dao2023flashattention2}
online softmax, recovering the full asymptotic gap on hardware.

\paragraph{Contributions.}
\begin{itemize}\itemsep0pt
  \item An ancestor-sparse Triton tree-attention kernel with
        $\bigO{\numtokens \!\cdot\! \depth}$ FMAs and HBM traffic
        (Sec.~\ref{sec:method}).
  \item Kernel-level wins of up to $7.75\times$ over FlashInfer and
        $16.99\times$ over DeFT in the deep/wide regime
        (Sec.~\ref{sec:eval-kernel}, Tables~\ref{tab:micro-flashinfer}--\ref{tab:micro-deft}).
  \item An honest end-to-end account: at short context our $3$--$4\times$
        kernel win at moderate tree size translates to only an
        $\approx\!2\%$ verification-cost reduction
        ($\approx\!1.04\times$ E2E) because attention's share is small
        in that regime (Sec.~\ref{sec:eval-honest}).
  \item Long-context end-to-end results showing
        $\bigO{\numtokens \!\cdot\! \depth}$ scaling in action: up to
        $1.61\times$ on EAGLE-3 at $\prefixlen{=}10\,$k and
        $1.40\times$ on Sequoia at $\numtokens{=}1024$
        (Secs.~\ref{sec:eval-prefix}--\ref{sec:eval-sequoia}).
  \item A reframing of the next bottleneck. With attention removed,
        wall time still rises strictly with tree size while accepted
        tokens per step plateau: the ceiling is acceptance-rate decay,
        a draft-model property, not a kernel property
        (Sec.~\ref{sec:eval-acceptance}).
\end{itemize}

\section{Background}
\label{sec:background}

\paragraph{Tree verification.}
Given a target model and a draft tree of $\numtokens$ nodes in BFS
order, the verifier runs one forward pass over all $\numtokens$ tokens.
The attention mask is \emph{ancestor-causal}: query at position $i$ may
attend to position $j$ iff $j$ is an ancestor of $i$ in the tree, or
$j {=} i$. For a complete $\branching$-ary tree of depth $\depth$,
$\parent(k){=}\lfloor(k{-}1)/\branching\rfloor$ for $k{>}0$ and
$\numtokens{=}(\branching^{\depth+1}{-}1)/(\branching{-}1)$.

\paragraph{The quadratic in current kernels.}
FlashInfer iterates over $\bigO{\numtokens}$ KV blocks per query and
applies a tree mask, costing $\bigO{\numtokens^2}$ per sequence.
DeFT~\cite{yao2024deft} restricts itself to strictly balanced trees and
factorises along the tree topology, but on the trees it supports it is
slower than FlashInfer in our measurements. Both kernels lock the
optimiser into a small-tree regime: not for algorithmic reasons but for
systems reasons.

\section{Method: \method{}}
\label{sec:method}

We replace the KV-block loop of a standard FlashAttention-2 kernel with
an \emph{ancestor walk} of exactly $\depth_{\max}{+}1$ iterations. Each
iteration loads $K[\parent^s(q)]$ and $V[\parent^s(q)]$ for ancestor
depth $s$ and feeds the result into the standard FA-2 online-softmax
update --- which is agnostic to whether the per-step contribution is a
$\textsc{Block\_N}$-block or a single ancestor score.
A duplicate guard $\text{prev}\!\neq\!\text{cur}$ turns iterations
beyond the root into no-ops, letting shallow queries share the same
kernel without branching.

A subtle PTX issue arises with non-power-of-two $\branching$: the
natural \verb|cur = (cur - 1) // b| compiles to an unsigned multiply-high,
turning $\text{cur}{=}0$ into a spurious ancestor index. We clamp
\emph{before} dividing: \verb|cur = max(cur - 1, 0) // b|.

\begin{algorithm}[t]
\caption{Ancestor-sparse FlashAttention-2 inner loop (single Q-tile, single head).}
\label{alg:kernel}
\begin{algorithmic}[1]
\Require Q, K, V; sequence offsets, branching factor $\branching$, max depth $D$
\State $\text{acc} \gets 0;\;\; m_i \gets -\infty;\;\; l_i \gets 0$
\State $\text{cur} \gets q\_\text{positions};\;\; \text{prev} \gets -1$
\For{$s = 0$ \textbf{to} $D$}                            \Comment{ancestor walk}
  \State $\text{is\_new} \gets (\text{cur} \neq \text{prev}) \wedge \text{valid}$
  \State $k_a \gets \mathrm{load}(K[\text{seq\_start} + \text{cur}])$
  \State $v_a \gets \mathrm{load}(V[\text{seq\_start} + \text{cur}])$
  \State $s \gets \langle q, k_a\rangle \cdot \tfrac{1}{\sqrt{D_h}}\;\;\text{(masked $-\infty$ where not is\_new)}$
  \State $m_{\text{new}} \gets \max(m_i, s);\;\; \alpha \gets e^{m_i - m_{\text{new}}};\;\; p \gets e^{s - m_{\text{new}}}$
  \State $\text{acc} \gets \alpha \cdot \text{acc} + p \cdot v_a;\;\; l_i \gets \alpha\cdot l_i + p;\;\; m_i \gets m_{\text{new}}$
  \State $\text{prev} \gets \text{cur};\;\; \text{cur} \gets \max(\text{cur}-1,0)\,\,/\!/\,\,\branching$
\EndFor
\State \Return $\text{acc} / l_i$
\end{algorithmic}
\end{algorithm}

\paragraph{Complexity.}
For a complete $\branching$-ary tree of depth $\depth$ with head
dimension $\headdim$ and $\nheads$ heads, the kernel performs
$\bigO{\batchsize\!\cdot\!\numtokens\!\cdot\!(\depth{+}1)\!\cdot\!\headdim\!\cdot\!\nheads}$
FMAs and streams the same order of K/V bytes, versus
$\bigO{\batchsize\!\cdot\!\numtokens^2\!\cdot\!\headdim\!\cdot\!\nheads}$
for FlashInfer.

\section{Evaluation}
\label{sec:eval}

\subsection{Setup}
\label{sec:eval-setup}

Kernel-level micro-benchmarks run on an RTX 4000 Ada (SM89, 20\,GB).
End-to-end runs use Llama-3.1-8B-Instruct with the EAGLE3-LLaMA3.1-8B
draft~\cite{li2025eagle3}, fp16 weights and activations; long-context
sweeps run on a single A100. Each micro-benchmark measurement is the
median of 50 CUDA-event-timed calls after 10 warmup iterations; both
kernels are timed identically. End-to-end results aggregate three
independent runs over a fixed prompt set; we report mean speedup over
the FlashInfer-backed pipeline (\emph{vanilla}) versus the same
pipeline with our kernel substituted in (\emph{ragged}). Prompts
that trigger early stopping (\textless 200 toks) are filtered out to avoid skewing the metrics.

\paragraph{Correctness.}
We compare \method{}'s per-token output tensors against a padded
PyTorch SDPA reference (fp32 math backend, exact tree mask) across
72 configurations spanning $\batchsize \in \{1,2,4,8\}$,
$\branching \in \{1,2,3\}$, $\depth \in \{1,\ldots,6\}$.
All 72 pass with atol\,=\,rtol\,=\,$10^{-2}$.
The global peak absolute error across all configurations and all
output elements is $\mathbf{1.95 \times 10^{-3}}$; the global mean
absolute error is $8.2 \times 10^{-8}$.
Both figures are consistent with fp16 online-softmax accumulation
rounding (fp16 unit-roundoff is $\approx 4.9 \times 10^{-4}$, and
accumulating $O(\depth)$ terms produces proportional drift).
No configuration produces a different top-1 logit under \method{}.

\begin{table}[!htbp]
\centering
\caption{Kernel speedup of \method{} over FlashInfer (left) and DeFT
(right) at $\batchsize{=}1$, $\prefixlen{=}0$ (RTX 4000 Ada, fp16).
Rows: tree depth $\depth$; columns: branching factor $\branching$.
Both baselines scale as $\bigO{\numtokens^2}$; \method{} scales as
$\bigO{\numtokens\!\cdot\!\depth}$, so speedups grow monotonically with
both axes. DeFT is strictly slower than FlashInfer at every
 configuration tested.}
\label{tab:micro-combined}
\vspace{4pt}
\begin{subtable}[t]{0.56\linewidth}
\centering
\caption{\method{} vs.\ FlashInfer}
\label{tab:micro-flashinfer}
\small
\begin{tabular}{l rrrrrrr}
\toprule
$\depth \backslash \branching$ & 4 & 8 & 12 & 16 & 20 & 24 & 28 \\
\midrule
28 & 1.59 & 2.20 & 2.86 & 5.82 & 6.32 & 6.65 & \textbf{7.75} \\
20 & 1.55 & 2.60 & 2.48 & 2.90 & 3.28 & 6.01 & 6.72 \\
14 & 0.93 & 1.80 & 2.69 & 2.58 & 2.75 & 2.81 & 3.27 \\
10 & 0.97 & 1.68 & 2.06 & 2.86 & 3.17 & 2.68 & 2.87 \\
 7 & 0.98 & 1.07 & 1.63 & 1.97 & 2.31 & 2.95 & 3.09 \\
 5 & 0.96 & 1.00 & 1.06 & 1.88 & 1.92 & 2.26 & 2.90 \\
 3 & 0.97 & 0.98 & 0.99 & 1.03 & 1.04 & 1.80 & 1.90 \\
\bottomrule
\end{tabular}
\end{subtable}%
\hfill
\begin{subtable}[t]{0.40\linewidth}
\centering
\caption{\method{} vs.\ DeFT (balanced trees)}
\label{tab:micro-deft}
\small
\begin{tabular}{l rrrr}
\toprule
$\depth \backslash \branching$ & 7 & 14 & 21 & 28 \\
\midrule
28 &  9.27 & 11.07 & 15.23 & \textbf{16.99} \\
21 &  6.74 & 10.26 & 14.00 & 16.86 \\
14 &  4.15 & 10.50 & 11.23 & 13.03 \\
 7 &  1.99 &  4.37 &  6.34 & 10.43 \\
\bottomrule
\end{tabular}
\end{subtable}
\end{table}
\subsection{Kernel-level: \method{} scales linearly}
\label{sec:eval-kernel}

Tables~\ref{tab:micro-flashinfer}--\ref{tab:micro-deft} sweep
$(\branching,\depth)$ at $\batchsize{=}1$ and no prefix length
to measure pure tree attention cost.

\paragraph{Near-parity at small trees (top-left corner).}
At $(\depth{\le}5,\branching{\le}8)$, both kernels report
speedups within $\approx\!0.97\times$ of each other.
This is expected: at small tree sizes the attention computation
itself is so brief (sub-microsecond) that kernel launch overhead,
cuDNN/Triton dispatch, and memory-subsystem latency dominate on
both sides equally. Neither kernel is limited by its FMA count;
both are limited by the cost of starting a CUDA kernel in the
first place. This is the regime production optimisers currently
target, and the result is reassuring: \method{} does not
regress the common case.

\paragraph{Super-linear growth in FlashInfer.}
As depth and branching factor increase, FlashInfer must iterate
over progressively more KV-cache blocks per query. In a
$\branching$-ary tree of depth $\depth$, each query's row of
the attention matrix is $\numtokens$-wide; FlashInfer processes
all $\numtokens$ key-value pairs and masks out the non-ancestors.
Both the FMA count and the HBM traffic therefore scale as
$\bigO{\numtokens^2}$ over the batch. \method{}, by contrast,
walks only the $\depth{+}1$ ancestor positions and processes
no others, so its cost scales as $\bigO{\numtokens\!\cdot\!(\depth{+}1)}$.
Since $\numtokens$ itself grows exponentially in $\depth$ for
fixed $\branching$ (and polynomially in $\branching$ for fixed
$\depth$), the ratio diverges rapidly in both axes --- exactly
what the tables show.

\paragraph{Peak gap at the deep/wide extreme.}
The deepest, widest configuration
$(\branching,\depth){=}(28,28)$ reaches $7.75\times$ over FlashInfer
and $16.99\times$ over DeFT. DeFT's larger gap arises because it uses a
different tree-factorisation strategy that, while asymptotically
similar to FlashInfer, incurs higher constant-factor overhead on
the GPU --- our measurements confirm DeFT is slower than
FlashInfer at every measured configuration
(Table~\ref{tab:micro-deft}).
Fig.~\ref{fig:kernel-scaling} plots latency directly against
$\numtokens$: the FlashInfer curves bend upward super-linearly
while \method{}'s curves remain nearly flat.


\begin{figure}[!htbp]
\centering
\includegraphics[width=0.65\linewidth]{fig3_kernel_latency_vs_tree_size.pdf}
\caption{Kernel latency vs.\ tree size for \method{} (solid) and
FlashInfer (dashed) at $\batchsize{=}1$ and $\batchsize{=}4$
($\prefixlen{=}0$). FlashInfer bends up super-linearly; \method{}
remains nearly flat once $\depth$ stops growing.}
\label{fig:kernel-scaling}
\end{figure}

\FloatBarrier

\subsection{End-to-end baseline with no added prefix}
\label{sec:eval-honest}

We state the unfavourable case first, because it is the
scientifically necessary control.

Fig.~\ref{fig:eagle-l0} sweeps EAGLE-3 tree size at $\prefixlen{=}0$
(no shared prompt prefix). Even though \method{} is $3$--$4\times$
faster than FlashInfer at $tt{\approx}330$ in kernel isolation, the
end-to-end gain collapses to only $\approx\!1.04\times$. The reason
is a straightforward application of Amdahl's law: at short context,
the verifier step is dominated by the MLP forward pass (which
\method{} does not touch), by draft-model inference, and by tree
construction and sampling overhead. The attention sub-step that
\method{} accelerates contributes a small fraction of total step
time. This is shown by the ~2\% relative drop in verify cost.

This is the expected consequence of optimising a non-bottleneck,
we report it strictly for ablation. The value of \method{}
at short context is the result shown in
Tables~\ref{tab:micro-flashinfer}--\ref{tab:micro-deft}: parity at current configs.

\begin{figure}[!htbp]
\centering
\includegraphics[width=0.55\linewidth]{fig6_eagle-3_speedup_vs_tree_size.pdf}
\caption{EAGLE-3 end-to-end throughput speedup and relative verify-cost
reduction at zero prefix length, sweeping total tokens per tree.
}
\label{fig:eagle-l0}
\end{figure}

\FloatBarrier

\subsection{The deeper problem: cost outpaces acceptance}
\label{sec:eval-acceptance}

With the kernel's asymptotic advantage established, we turn to a
fundamental constraint that operates entirely above the attention
layer. Fig.~\ref{fig:accept-vs-cost} plots two quantities from
EAGLE-3 as a function of tree size, both measured at $\prefixlen{=}0$
and kernel-agnostically (i.e., the effect holds regardless of which
attention kernel is used).

\paragraph{Why wall time rises.}
The target model runs one forward pass over all $\numtokens$ tree
tokens. The MLP is the dominant cost and scales linearly in
$\numtokens$; KV-cache prefill grows with the shared prefix; and
layer normalisation, residuals, and scheduling overhead all add
constant multiples of $\numtokens$. Even with a linear-time attention
kernel, the step's wall time is therefore a strictly increasing
function of $\numtokens$. There is no free lunch: more candidates
always cost more GPU time.

\paragraph{Why acceptance rate per token collapses.}
The draft model predicts token $k$ at depth $s$ of the tree
conditioned on the $s$ ancestor tokens it was asked to extend. As $s$
grows, the draft's hidden states must capture an increasingly long
hypothetical continuation that the target model has never actually
sampled --- the draft is predicting the target's predictions of its
own future predictions. The EAGLE-3~\cite{li2025eagle3} ablations
show this saturation directly; recent work on hidden-state
drift~\cite{drift2026} formalises it: representational drift between
the draft and target models accumulates with candidate-horizon depth,
capping per-token acceptance probability independent of how much
verifier compute is available.

\begin{figure}[!htbp]
\centering
\begin{subfigure}[t]{0.48\linewidth}\centering
\includegraphics[width=\linewidth]{fig4_acceptance_rate_vs_cost.pdf}
\subcaption{Per-token acceptance rate (left axis) and verifier wall
time (right axis) vs.\ tree size. Wall time rises near-linearly
(MLP-dominated); acceptance rate falls monotonically due to
representational drift at deeper candidate horizons.}
\label{fig:accept-vs-cost}
\end{subfigure}\hfill
\begin{subfigure}[t]{0.48\linewidth}\centering
\includegraphics[width=\linewidth]{fig5_accepted_tokens_vs_cost.pdf}
\subcaption{Accepted tokens per step (left axis) and wall time (right
axis). Accepted tokens plateau rapidly while wall time grows linearly, 
the marginal tree node costs more than it returns in accepted
tokens beyond the smallest trees. The crossover occurs around
$\numtokens \approx 50$.}
\label{fig:accept-tokens-vs-cost}
\end{subfigure}
\caption{The acceptance--cost trade-off in EAGLE-3, kernel-agnostic
($\prefixlen{=}0$). \emph{Left}: the two opposing trends in raw form
--- wall time up, per-token acceptance rate down.
\emph{Right}: their product normalised to accepted tokens per step,
the quantity that directly determines throughput. Neither effect is a
kernel artifact: both persist identically under Vanilla and
\method{}. }
\label{fig:problem}
\end{figure}

The implication is the central call to action of this paper:
\emph{better draft models that maintain acceptance rates at scale
are the missing piece.} Quadratic attention has masked this
bottleneck for years by making deep trees impractical on systems
grounds alone. Now that the systems constraint is lifted, the
statistical constraint --- draft-model fidelity at depth --- stands
exposed. The design space that kernel limitations had closed is
once again open on its empirical merits.

\FloatBarrier

\subsection{\method{} removes attention as the excuse: prefix scaling}
\label{sec:eval-prefix}

Section~\ref{sec:eval-honest} showed that at short context, attention
is a small fraction of verifier wall time, so even a
multi-$\times$ kernel improvement leaves end-to-end throughput
nearly unchanged. The natural question is: what happens when attention's
share of step time actually grows?

Fig.~\ref{fig:ctx-sweep} answers this by fixing the tree shape at
$(\branching,\depth){=}(20,20)$ ($tt{=}343$) and sweeping prefix length
from $2048$ to $10240$ tokens on an A100. As $\prefixlen$ grows,
the KV-cache that each attention query must scan grows proportionally,
so FlashInfer's $\bigO{\numtokens^2}$ inner loop accumulates more
unused work. \method{}'s ancestor walk, by contrast, touches only
the $\depth{+}1{=}21$ ancestor positions per query regardless of
how long the KV cache is --- so its verify time is nearly
constant across the sweep. The consequence, shown in
Table~\ref{tab:e2e-prefix}, is that verify-cost reduction grows
monotonically: $-10.8\%$ at $\prefixlen{=}2048$ climbing to
$-28.0\%$ at $\prefixlen{=}10240$. End-to-end throughput inherits
this gain and reaches $\boldsymbol{1.61\times}$ at $\prefixlen{=}10240$ tokens.

\begin{table}[!htbp]
\centering
\caption{End-to-end throughput speedup of \method{} over FlashInfer on
EAGLE-3 with prefix sweep ($\branching{=}20,\depth{=}20,tt{=}343$,
A100). Verify cost reduction grows monotonically with $\prefixlen$.}
\label{tab:e2e-prefix}
\small
\begin{tabular}{l rrrrr}
\toprule
$\prefixlen$           & 2048 & 4096 & 6144 & 8192 & 10240 \\
\midrule
Verify ms (FlashInfer) & 0.547 & 0.613 & 0.648 & 0.672 & 0.683 \\
Verify ms (\method{})  & 0.488 & 0.492 & 0.495 & 0.492 & 0.492 \\
$\Delta$ verify cost   & $-10.8\%$ & $-19.7\%$ & $-23.6\%$ & $-26.7\%$ & $-28.0\%$ \\
E2E speedup            & $1.12\times$ & $1.31\times$ & $1.45\times$ & $1.57\times$ & $\boldsymbol{1.61\times}$ \\
\bottomrule
\end{tabular}
\end{table}

\begin{figure}[!htbp]
\centering
\includegraphics[width=0.60\linewidth]{fig7_a100_ctx_sweep.pdf}
\caption{End-to-end throughput speedup and verify-cost reduction vs.\
prefix length at fixed $(\branching,\depth){=}(20,20)$, $tt{=}343$
(A100, Llama-3.1-8B, EAGLE-3). FlashInfer's verify time grows with
$\prefixlen$ because each query scans the full KV cache;
\method{}'s verify time is nearly flat because each query walks
only its $\depth{+}1$ ancestors. The gap between the two verify
times translates directly into monotonically rising E2E throughput,
reaching $1.61\times$ at $\prefixlen{=}10\,$k.}
\label{fig:ctx-sweep}
\end{figure}

\FloatBarrier

Fig.~\ref{fig:n-sweep} approaches the same linear-vs-quadratic
gap from the orthogonal axis: it fixes $\prefixlen{=}4096$ and
sweeps tree size at the same A100 setup. At small trees,
attention's share of step time is small even at this context length,
and the two kernels are near-parity. As $(\branching,\depth)$ grows,
FlashInfer's inner loop must process an $\bigO{\numtokens^2}$
KV-scan for each additional query, while \method{} adds only
$\bigO{\depth}$ ancestor loads per query. The speedup therefore
grows with tree size and reaches $1.41\times$ at
$(\branching,\depth){=}(24,24)$ ($tt{=}494$). Together, the two
panels confirm a single underlying claim: the $\bigO{\numtokens\!\cdot\!\depth}$
vs.\ $\bigO{\numtokens^2}$ asymptotic gap is not an abstract kernel
measurement --- it plays out at the system level wherever and
whenever attention's share of the verifier step is non-trivial.

\begin{figure}[!htbp]
\centering
\includegraphics[width=0.60\linewidth]{fig8_a100_n_sweep.pdf}
\caption{End-to-end throughput speedup vs.\ tree size at fixed
$\prefixlen{=}4096$ (A100, Llama-3.1-8B, EAGLE-3). At small trees,
both kernels are near-parity because attention's share of step
time is small. As tree size grows, FlashInfer's $\bigO{\numtokens^2}$
inner loop begins to dominate its verify latency, and the speedup
grows monotonically, reaching $1.41\times$ at
$(\branching,\depth){=}(24,24)$, $tt{=}494$.}
\label{fig:n-sweep}
\end{figure}

\FloatBarrier

\subsection{\texorpdfstring{Sequoia: scaling trees up to $\numtokens{=}1024$}{Sequoia: scaling trees up to N=1024}}
\label{sec:eval-sequoia}

Sequoia~\cite{chen2024sequoia} is designed explicitly to scale trees
beyond what current systems treat as practical, making it the cleanest
available end-to-end test of the asymptotic distinction between
$\bigO{\numtokens \!\cdot\! \depth}$ and $\bigO{\numtokens^2}$.
Unlike EAGLE-3's dynamic-draft trees, Sequoia uses fixed
shapes --- the only variable across runs is $\numtokens$.

Table~\ref{tab:sequoia} and Fig.~\ref{fig:sequoia} show throughput
for both kernels across the size grid from $\numtokens{=}4$ to
$\numtokens{=}1024$. Through $\numtokens{=}256$, both kernels are
near-parity: the MLP forward pass dominates step time and FlashInfer's
quadratic inner loop has not yet grown large enough to stand out against
that background. The crossover becomes visible around $\numtokens{=}512$,
where FlashInfer's throughput drops sharply ($27.1$ tok/s) while
\method{} remains higher ($29.3$ tok/s, a $1.08\times$ advantage).
By $\numtokens{=}1024$, FlashInfer's attention cost has grown to the
point that it visibly throttles the pipeline: throughput falls to
$11.5$ tok/s, less than a quarter of the peak near $\numtokens{=}64$.
\method{}'s throughput at $\numtokens{=}1024$ is $16.2$ tok/s ---
a $\boldsymbol{1.40\times}$ gain --- precisely because its attention
term has not grown with $\numtokens^2$.

Stochastic sampling introduces a small per-step divergence of
$\pm 0.08$ accepted tokens between the two kernels, visible on
Fig.~\ref{fig:sequoia}'s right axis. This is the expected
variability from temperature sampling, not a correctness issue:
per the parity check in Sec.~\ref{sec:eval-setup}, the two kernels
produce identical logits.

\begin{table}[!htbp]
\centering
\caption{Sequoia tree-size sweep, end-to-end throughput
(tokens / second) and \method{} speedup. Sequoia fixes shape
and varies only $\numtokens$, so this isolates the
linear-vs-quadratic effect.}
\label{tab:sequoia}
\small
\begin{tabular}{r rr r}
\toprule
$\numtokens$ & FlashInfer (tok/s) & \method{} (tok/s) & E2E speedup \\
\midrule
   4 & 48.5  & 48.5  & $1.00\times$ \\
   8 & 55.3  & 55.2  & $1.00\times$ \\
  16 & 62.8  & 58.5  & $0.93\times$ \\
  32 & 63.5  & 62.4  & $0.98\times$ \\
  64 & 64.2  & 63.6  & $0.99\times$ \\
 128 & 59.6  & 59.8  & $1.00\times$ \\
 256 & 48.7  & 49.1  & $1.01\times$ \\
 512 & 27.1  & 29.3  & $1.08\times$ \\
1024 & 11.5  & 16.2  & $\boldsymbol{1.40\times}$ \\
\bottomrule
\end{tabular}
\end{table}

\begin{figure}[!htbp]
\centering
\includegraphics[width=0.55\linewidth]{fig9_Sequoia_speedup.pdf}
\caption{Sequoia~\cite{chen2024sequoia} end-to-end throughput speedup
across the size grid. Sequoia uses fixed trees and exploits
larger sizes directly. Stochastic sampling introduces mild divergence
in accepted tokens per step between kernels (right axis, $\pm 0.08$),
but the throughput trend is clear: $1.08\times$ at $\numtokens{=}512$
and $1.40\times$ at $\numtokens{=}1024$.}
\label{fig:sequoia}
\end{figure}

\FloatBarrier

\subsection{The remaining frontier}
\label{sec:eval-frontier}

Combining the threads:
(i)~\method{} converts $\bigO{\numtokens^2}$ tree-attention into
$\bigO{\numtokens \!\cdot\! \depth}$ at the kernel level
(Sec.~\ref{sec:eval-kernel}).
(ii)~Where attention's share of verifier time is large --- long context
or large tree --- the kernel translates directly into end-to-end speedup
of $1.4$--$1.6\times$ (Secs.~\ref{sec:eval-prefix}--\ref{sec:eval-sequoia}).
(iii)~Where attention's share is small --- short context, small tree
--- the kernel is honestly near-parity, by construction
(Sec.~\ref{sec:eval-honest}).
(iv)~Across all regimes, accepted tokens per step plateau as trees
grow, so wall-time savings from larger trees do not translate into
proportionally more accepted tokens (Sec.~\ref{sec:eval-acceptance}).

The first three are kernel results. The fourth is not.
With attention removed as the bottleneck, the next frontier is
\emph{acceptance-rate decay}: a property of the draft model and its
relationship to the target model, not of the verifier kernel. We hand
the design space back to draft-model researchers --- tree-construction
policies (e.g., Sequoia's scaling acceptance curves) can now be
revisited on their statistical merits rather than against a quadratic
verifier.

\section{Related Work}
\label{sec:related}

\paragraph{Speculative decoding.}
Speculative sampling~\cite{leviathan2023fast,chen2023accelerating} and
its tree-based descendants~\cite{cai2024medusa,li2024eagle,li2024eagle2,
li2025eagle3,chen2024sequoia,miao2024specinfer,spector2023accelerating}
have become a standard inference acceleration technique. Tree-construction
policies in this line target tree shape \emph{given} a fixed attention
kernel. \method{} is complementary: it changes the kernel, expanding the
design space the optimiser operates in.

\paragraph{Acceptance-rate ceilings.}
EAGLE-3~\cite{li2025eagle3} reports diminishing returns as draft trees
grow and attributes this to the draft's representational fidelity at
long horizons. Concurrent work on hidden-state
drift~\cite{drift2026} formalises the same effect: the
draft's hidden states drift from the target's as the candidate horizon
grows, capping acceptance regardless of how much verifier compute is
available. Our findings are consistent with theirs and locate the
remaining bottleneck on the draft side.

\paragraph{Attention kernels.}
FlashAttention~\cite{dao2022flashattention,dao2023flashattention2}
established the IO-aware online-softmax recipe we extend.
FlashInfer~\cite{ye2025flashinfer} is the de-facto production kernel
and our primary baseline; DeFT~\cite{yao2024deft} is slower than FlashInfer in our measurements across the board
(Table~\ref{tab:micro-deft}).
Block-sparse and mask-aware kernels~\cite{xformers2022,gao2024seerattention}
exploit structured sparsity at $\bigO{\text{block}}$ granularity rather
than the $\bigO{\depth}$ ancestor structure specific to speculation
trees.

\section{Conclusion}
\label{sec:conclusion}

We presented \method{}, an ancestor-sparse FlashAttention-2 kernel for
tree-structured speculative decoding with
$\bigO{\numtokens \!\cdot\! \depth}$ work and HBM traffic in place of
the $\bigO{\numtokens^2}$ inner loop of production tree-attention
kernels. Kernel-level it reaches $7.75\times$ over FlashInfer and
$16.99\times$ over DeFT in the deep/wide regime; end-to-end it
delivers $1.61\times$ on EAGLE-3 at long context and $1.40\times$ on
Sequoia at $\numtokens{=}1024$. We are equally explicit about where it
does not help: at short context, attention is a small share of step
time, and the $\approx\!1.04\times$ honest baseline reflects that.
Across all regimes, accepted tokens per step plateau as trees grow ---
a draft-model property our kernel cannot improve. With attention
removed as the systematic blocker, draft-model fidelity is the next
frontier; the design space the field had closed for systems reasons is
once again open on its statistical merits.

\paragraph{Limitations.}
Our empirical evaluation restricts the target model (Llama-3.1-8B) and
hardware (RTX 4000 Ada and A100) for ablation and comparison.
While the $\bigO{\numtokens\!\cdot\!\depth}$ algorithm itself relies purely on
the causal independence of tree branches and is strictly model-agnostic,
constants in our Triton implementation may shift under different GPU
architectures (e.g., AMD instinct, Nvidia H100) or under extreme
head-dimension scaling, which we leave to future profiling.

All code, experiments and full results are available at \url{https://github.com/saifmb0/QAttention}.

\bibliographystyle{plain}
\begin{thebibliography}{99}\small
\bibitem{leviathan2023fast} Y.\ Leviathan, M.\ Kalman, Y.\ Matias.
  Fast Inference from Transformers via Speculative Decoding. \emph{ICML}, 2023.
\bibitem{chen2023accelerating} C.\ Chen et al.
  Accelerating Large Language Model Decoding with Speculative Sampling.
  \emph{arXiv:2302.01318}, 2023.
\bibitem{cai2024medusa} T.\ Cai et al.
  Medusa: Simple LLM Inference Acceleration Framework with Multiple Decoding Heads.
  \emph{ICML}, 2024.
\bibitem{li2024eagle} Y.\ Li et al.
  EAGLE: Speculative Sampling Requires Rethinking Feature Uncertainty. \emph{ICML}, 2024.
\bibitem{li2024eagle2} Y.\ Li et al.
  EAGLE-2: Faster Inference of Language Models with Dynamic Draft Trees.
  \emph{EMNLP}, 2024.
\bibitem{li2025eagle3} Y.\ Li et al.
  EAGLE-3: Scaling up Inference Acceleration of Large Language Models via Training-Time Test.
  \emph{arXiv:2503.01840}, 2025.
\bibitem{chen2024sequoia} Z.\ Chen et al.
  Sequoia: Scalable, Robust, and Hardware-aware Speculative Decoding.
  \emph{NeurIPS}, 2024.
\bibitem{miao2024specinfer} X.\ Miao et al.
  SpecInfer: Accelerating Generative LLM Serving with Tree-based Speculative Inference.
  \emph{ASPLOS}, 2024.
\bibitem{spector2023accelerating} B.\ Spector, C.\ Re.
  Accelerating LLM Inference with Staged Speculative Decoding.
  \emph{arXiv:2308.04623}, 2023.
\bibitem{dao2022flashattention} T.\ Dao et al.
  FlashAttention: Fast and Memory-Efficient Exact Attention with IO-Awareness.
  \emph{NeurIPS}, 2022.
\bibitem{dao2023flashattention2} T.\ Dao.
  FlashAttention-2: Faster Attention with Better Parallelism and Work Partitioning.
  \emph{ICLR}, 2024.
\bibitem{ye2025flashinfer} Z.\ Ye et al.
  FlashInfer: Efficient and Customizable Attention Engine for LLM Inference Serving.
  \emph{MLSys}, 2025.
\bibitem{yao2024deft} J.\ Yao et al.
  DeFT: Decoding with Flash Tree-Attention for Efficient Tree-Search-Based LLM Inference.
  \emph{ICLR}, 2025.
\bibitem{drift2026} Qwen.
  When Hidden States Drift: Acceptance-Rate Decay in Tree-Based Speculative Decoding.
  \emph{arXiv:2604.26412}, 2026.
\bibitem{xformers2022} B.\ Lefaudeux et al.
  xFormers: A modular and hackable Transformer modelling library.
  \emph{Meta AI}, 2022.
\bibitem{gao2024seerattention} Y.\ Gao et al.
  SeerAttention: Learning Intrinsic Sparse Attention in Your LLMs.
  \emph{arXiv:2410.13276}, 2024.
\end{thebibliography}
\end{document}
